\section{Evaluation Metrics and Statistical Protocol}
\label{app:evaluation_metrics}

This appendix defines the evaluation protocol for the current six-mode C--K representation. It replaces the earlier macro-TED formulation, which used an obsolete $P,K,V,C$ partition and a raw edit count. Here, $V$ is not a generated C--K state. It denotes value meanings expressed in the generated or human response.

Let $\mathcal U_s$ be the set of participants with a valid prediction and gold response for segment $s$, and let $N_s=|\mathcal U_s|$. Metrics with per-participant scores are reported as participant means; cohort-level diversity metrics are defined below. We use the frozen BGE-M3 dense encoder $e(\cdot)$ (1024 dimensions, L2-normalized embeddings) for every embedding-based metric.

\subsection{Seg2 Trajectory Fidelity}

\paragraph{KC Normalized Edit Similarity (KC-NES).}
For participant $u$, let $y_u$ be the gold C--K sequence after removing auxiliary $V$ labels while preserving the order of the remaining operations. Let $\hat y_u$ be the predicted C--K sequence after removing boundary controls such as \textsc{Stop}. We define
\begin{equation}
\operatorname{KC\text{-}NES}_u
=1-\frac{\operatorname{Lev}(\hat y_u,y_u)}
{\max(|\hat y_u|,|y_u|,1)},
\label{eq:kc_nes_individual}
\end{equation}
where insertion, deletion, and substitution each have unit cost. The reported score is
\begin{equation}
\overline{\operatorname{KC\text{-}NES}}
=\frac{1}{N_{\mathrm{Seg2}}}\sum_{u\in\mathcal U_{\mathrm{Seg2}}}
\operatorname{KC\text{-}NES}_u.
\label{eq:kc_nes}
\end{equation}
KC-NES ranges from 0 to 1 and jointly penalizes incorrect operations, order, and sequence length.

\paragraph{KC Own-User Margin and AUC.}
KC-NES alone does not establish personalization because a cohort-modal trajectory may achieve moderate similarity for many participants. We therefore compare each prediction against every participant's gold trajectory. Let
\begin{equation}
q^{KC}_{uv}=\operatorname{KC\text{-}NES}(\hat y_u,y_v).
\end{equation}
The participant-level own-user margin is
\begin{equation}
m_u^{KC}=q^{KC}_{uu}-\frac{1}{N_{\mathrm{Seg2}}-1}
\sum_{v\ne u}q^{KC}_{uv}.
\label{eq:kc_own_margin}
\end{equation}
Because short discrete trajectories produce many ties, our primary trajectory-personalization score is a tie-aware pairwise AUC:
\begin{equation}
\operatorname{AUC}^{KC}_u=\frac{1}{N_{\mathrm{Seg2}}-1}
\sum_{v\ne u}\left[
\mathbb I(q^{KC}_{uu}>q^{KC}_{uv})+
\frac12\mathbb I(q^{KC}_{uu}=q^{KC}_{uv})\right].
\label{eq:kc_own_auc}
\end{equation}
An AUC of 0.5 indicates no identity-discriminative trajectory signal; a margin of zero indicates no absolute own-user advantage. Invalid predictions are treated as empty trajectories, yielding margin 0 and AUC 0.5 against nonempty gold trajectories, while KC-NES remains 0.

\paragraph{Chain Diversity.}
To diagnose trajectory collapse, let $\hat y_u$ be the complete predicted sequence and let $\mathcal S$ be the set of unique predicted sequences. For each $a\in\mathcal S$, let $n_a$ denote its frequency:
\begin{equation}
n_a=\sum_{u\in\mathcal U_{\mathrm{Seg2}}}\mathbb I(\hat y_u=a).
\label{eq:sequence_frequency}
\end{equation}
The modal sequence share and chain diversity are then
\begin{equation}
\operatorname{TopShare}=\frac{\max_{a\in\mathcal S}n_a}{N_{\mathrm{Seg2}}}.
\label{eq:top_share}
\end{equation}
\begin{equation}
\operatorname{ChainDiversity}=1-\operatorname{TopShare}.
\label{eq:chain_diversity}
\end{equation}
This metric diagnoses collapse only. It is not evidence that a trajectory matches the intended participant. We additionally report $|\mathcal S|$, the number of unique predicted sequences.

\paragraph{Pairwise KC Diversity Retention.}
Exact-chain counts do not detect near-identical trajectories. We therefore define
\begin{equation}
D_{pred}=\frac{2}{N(N-1)}\sum_{u<v}
\left[1-\operatorname{KC\text{-}NES}(\hat y_u,\hat y_v)\right],
\end{equation}
and define $D_{gold}$ analogously. Pairwise KC diversity retention is
\begin{equation}
\operatorname{KCDR}=D_{pred}/D_{gold}.
\end{equation}
We report both component distances. KCDR=0 denotes complete trajectory collapse and KCDR=1 denotes the same aggregate dispersion as human gold; values above one indicate over-dispersion rather than necessarily better predictions.

\paragraph{Trajectory Geometry Correlation.}
To test whether the system preserves who is similar to whom rather than merely producing varied chains, we compute
\begin{equation}
\operatorname{TGC}=\rho_{\mathrm{Spearman}}
\left(\{1-\operatorname{NES}(\hat y_u,\hat y_v)\}_{u<v},
\{1-\operatorname{NES}(y_u,y_v)\}_{u<v}\right),
\end{equation}
using average ranks for ties. We define TGC as zero when all predicted pairwise distances are constant. TGC is a cohort-level diagnostic; participant pairs are not treated as independent observations.

\subsection{Textual Own-User Fidelity}

\paragraph{BGE Own-User Margin.}
For a generated response $\hat Y_u$ and a gold response $Y_v$, define their pairwise semantic similarity as
\begin{equation}
q_{uv}=\cos(e(\hat Y_u),e(Y_v)).
\label{eq:bge_pairwise_similarity}
\end{equation}
The own-user and other-user similarities are
\begin{equation}
s_u^{\mathrm{own}}=q_{uu}.
\label{eq:bge_own_similarity}
\end{equation}
\begin{equation}
s_u^{\mathrm{other}}=
\frac{1}{N_s-1}\sum_{v\in\mathcal U_s\setminus\{u\}}q_{uv}.
\label{eq:bge_other_similarity}
\end{equation}
The per-participant and system-level margins are
\begin{equation}
m_u=s_u^{\mathrm{own}}-s_u^{\mathrm{other}},
\qquad
\overline m_s=\frac{1}{N_s}\sum_{u\in\mathcal U_s}m_u.
\label{eq:bge_margin}
\end{equation}
A positive margin means that the response is closer to its own gold response than to the average gold response of other participants. We compute this metric separately for Seg1, Seg2, and Seg3; these segments are not pooled.
For interpretability, we also report the means of $s_u^{\mathrm{own}}$ and $s_u^{\mathrm{other}}$ alongside the margin.

\subsection{Seg2 Value Alignment}

\paragraph{Value Semantic Similarity.}
Independent blind annotation extracts value-bearing spans from the generated and gold responses. Let $\mathcal V_u^{\mathrm{pred}}$ and $\mathcal V_u^{\mathrm{gold}}$ be the respective spans concatenated in original order. Conditional on both being nonempty,
\begin{equation}
S_u^V=\cos\left(e(\mathcal V_u^{\mathrm{pred}}),
e(\mathcal V_u^{\mathrm{gold}})\right).
\label{eq:value_semantic_similarity}
\end{equation}
We additionally report value-presence precision, recall, and F1; the four counts for empty/nonempty generated and gold value spans; and the effective sample size of the conditional score. Responses with both spans empty are not assigned similarity 1.

Because the conditional cohort can differ across systems, we use a coverage-adjusted score for system comparison on the fixed subset with nonempty gold value spans:
\begin{equation}
\operatorname{CAS}_u^V=
\mathbb I[\mathcal V_u^{\mathrm{pred}}\ne\varnothing]
\max(0,S_u^V).
\label{eq:coverage_adjusted_value}
\end{equation}

\paragraph{V30 Matched-Minus-Mismatched Fit.}
For each nonempty generated value span, two blinded raters score its fit to the participant's matched V30 profile and to three fixed, no-fixed-point mismatched profiles from the same outer-test fold. Let $r_u^{\mathrm{matched}}$ be the matched rating and $r_{uk}^{\mathrm{mismatched}}$ the $k$-th mismatched rating. The participant-level contrast is
\begin{equation}
\Delta_u^{V30}
=r_u^{\mathrm{matched}}
-\frac{1}{3}\sum_{k=1}^{3}r_{uk}^{\mathrm{mismatched}}.
\label{eq:v30_matched_margin}
\end{equation}
We report the participant mean, a participant-bootstrap confidence interval, and the effective sample size. Positive values indicate greater compatibility with the target profile than with control profiles. The primary analysis leaves responses without a generated value span as missing; a coverage-adjusted sensitivity analysis assigns their contrast to zero on the full fixed cohort.

\subsection{Seg1 and Seg3 Human Alignment}

\paragraph{Human Intuition Match.}
For Seg1, two blinded raters score, on a five-point scale, the agreement between a generated response and the target gold response in initial reaction, central concern, and attitude direction. For participant $u$, the score is
\begin{equation}
I_u=\frac{r_{u1}+r_{u2}}{2}.
\label{eq:intuition_match}
\end{equation}
We report the participant mean and confidence interval.

\paragraph{Human Criterion Match.}
For Seg3, raters use the same five-point scale to assess agreement in the final boundary, acceptable compromise, and value direction. Scores are aggregated as in Equation~\ref{eq:intuition_match}.

\paragraph{Criterion Uniqueness.}
Let $C_u$ denote the generated Seg3 criterion or bottom-line statement, and let $d_{uv}=\cos(e(C_u),e(C_v))$ be the similarity between two participants' criteria. Its between-participant diversity is
\begin{equation}
\mathrm{CU}
=1-\frac{2}{N_{\mathrm{Seg3}}(N_{\mathrm{Seg3}}-1)}
\sum_{u<v}d_{uv},
\label{eq:criterion_uniqueness}
\end{equation}
where $\mathrm{CU}$ denotes Criterion Uniqueness and the sum ranges over unordered pairs in $\mathcal U_{\mathrm{Seg3}}$.
High uniqueness only indicates that generated criteria differ. It does not establish that they match the intended participants.

\subsection{Statistical Reporting}

The participant is the sole unit of inference. For system comparisons, we apply a participant-level paired bootstrap with 10,000 resamples and report percentile 95\% confidence intervals for the paired difference. When multiple random seeds are used, their outcomes are first aggregated within participant. Human ratings additionally report quadratic-weighted Cohen's $\kappa$, ICC(2,1), and the effective number of participants. Value-semantic results always report coverage and effective sample size; a conditional mean alone is not used for cross-system conclusions.
